\section{Related Work}
\label{sec:related}

\paragraph{Scaling laws for compression.} Existing compression laws
differ in what they take as input and what they predict. Sparsity
enters as a multiplicative factor on the size term of a pretraining
loss law \citep{frantar2023scaling}; precision enters through an
effective-parameter count with a post-training degradation term that
grows with pretraining tokens \citep{kumar2024scaling}; formats used
during training collapse onto one capacity scalar
\citep{panferov2025unified}; the P$^2$ law predicts the post-training
loss of a pruned model from its size and loss before pruning, the
pruning rate, and the post-training tokens, validated on held-out
tokens, sizes, and rates \citep{zhang2024scaling}; and a distilled
student follows a broken power law in teacher loss during
pretraining-scale logit distillation \citep{busbridge2025distillation}.
Their targets are aggregate losses; we follow the P$^2$ law in
admitting the pre-compression loss and add a per-capability endpoint
measured the same way for three methods, a development panel in which
size and stage vary separately, and held-out tests with same-input
baselines.

\paragraph{Task-conditioned evaluation of compression.} Pruning
disproportionately harms knowledge-intensive tasks
\citep{jaiswal2023compressing}, low-bit quantization disproportionately
harms reasoning \citep{liu2025quantization}, and selective truncation
can improve specific behaviors \citep{sharma2023laser}. Laws with
task-level targets exist: \citet{sengupta2025compression} fit pruning
laws to test loss and averaged zero-shot accuracy against the pruning
ratio and recovery tokens, with per-task exponents and transfer to
unseen models in their revised version; \citet{zhou2025task} stratify
post-training quantization laws into memorization, application, and
reasoning strata for one quantizer; task-specific distillation has
been mapped over dataset size, compression ratio, and supervision
format \citep{ghita2026task}. Our relations use a
continuous conditional loss, admit the dense capability anchor and
training tokens as inputs, and treat the three methods with one
protocol.

\paragraph{Fitting and validating empirical laws.} Downstream accuracy
is predictable mainly through continuous log-likelihood measurements
\citep{schaeffer2023emergent,schaeffer2024predicting,gadre2024language,bhagia2024establishing,lourie2025scaling},
which motivates our loss endpoint and its stated boundaries; fitted
scaling-law forms are fragile to the fitting protocol
\citep{li2025misfitting,besiroglu2024chinchilla}, and repeated tokens
are not worth new tokens in data-constrained pretraining
\citep{muennighoff2023scaling}, an effect our reuse-count coordinate
measures in the distillation setting. Gradient-based importance
underlies modern pruning
\citep{lecun1989optimal,frantar2023sparsegpt,sun2023simple}; we use
gradient statistics only as mechanism evidence
(Appendix~\ref{sec:app_mech}).
